Quan mesurem els valors de potencial elèctric en una cubeta obtenim la distribució representada a la figura, en què podem observar dues càrregues ($Q_1$ i $Q_2$), una de positiva i una de negativa.

\figura[0.55]{fig1.png}

\begin{apartat}{1,25}
Determineu de manera raonada quina és la càrrega positiva i quina la negativa. Segons la vostra resposta, dibuixeu la direcció i el sentit del camp elèctric al punt $A$.
\end{apartat}

\begin{apartat}{1,25}
Suposeu que un electró es mou del punt $A$ al punt $B$. Calculeu el treball que fa el camp elèctric durant aquest moviment. Quin treball fa el camp elèctric quan l'electró es mou del punt $A$ al punt $C$ passant per $B$?
\end{apartat}

\dades{%
  $|e| = 1,602\times 10^{-19}\ \mathrm{C}$.}
